% ===========================================================================
% Section 4: Empirical Strategy
%
% Ported from paper/paper_planned.pdf. The coauthor draft is partially
% written and partially scaffolded — estimating equation, first-stage
% equation, and validation-package table plans are final; several
% subsections contain author-written TODO markers flagging where
% analysis output is needed before the prose can be written.
%
% All coauthor scaffolding is preserved verbatim as LaTeX comments so
% no TODO is lost in the port. The analysis outputs that unblock the
% TODOs come from:
%   - First-stage F-stats: Phase 1/2 analysis (already in panel)
%   - Pre-period balance: Table 6 script (not yet written)
%   - Pre-trends: Table 7 script (not yet written)
%
% CITATION KEYS introduced:
%   (none beyond Section 3; equation structure reuses MA formalism)
% ===========================================================================

\section{Empirical Strategy}
\label{sec:identification}

\subsection{Estimating Equation}

Our main specification relates changes in district-level outcomes to
changes in market access:
\begin{equation}
\Delta Y_i = \beta \cdot \Delta \ln \mathrm{MA}^{\mathrm{full}}_i
             + X'_i \gamma + \varepsilon_i,
\label{eq:main}
\end{equation}
where $\Delta Y_i$ is the change in outcome, $\Delta \ln
\mathrm{MA}^{\mathrm{full}}_i$ is the change in log market access, and
$X_i$ includes baseline log market access, baseline log population, and
geographic characteristics. The parameter $\beta$ is the elasticity of
regional outcomes with respect to market access.

\subsection{Endogeneity Concerns}

The regressor $\Delta \ln \mathrm{MA}^{\mathrm{full}}_i$ is not
exogenous. The infrastructure changes that generate it --- rail
closures and road construction between 1960 and 1986 --- were policy
choices, not random assignments, and those choices plausibly responded
to the same district conditions that drive the outcomes.

The direction of the resulting bias is ambiguous because
infrastructure was allocated under two opposing logics. Under an
efficiency logic, investment follows expected growth and divestment
follows expected decline: roads are built, and rail retained, where
outcomes were already rising, while rail is closed where outcomes were
already falling. This aligns market-access gains with positive outcome
shocks and losses with negative shocks --- a positive correlation
between $\Delta \ln \mathrm{MA}^{\mathrm{full}}_i$ and $\varepsilon_i$
that biases the ordinary-least-squares estimate of $\beta$ upward.
Under a compensatory logic, investment instead targets lagging
districts --- for example, roads extended to areas that were declining
--- aligning market-access gains with negative outcome shocks and
biasing $\beta$ downward. Because both logics plausibly operated over
1960--1986, the sign of the net bias is ambiguous. We address the
endogeneity with two instruments that move market access through
variation plausibly orthogonal to district-level outcome shocks,
described next.

\subsection{Instrumental Variables}

We instrument $\Delta \ln \mathrm{MA}^{\mathrm{full}}_i$ with two
changes in market access, each computed on a counterfactual network:
one that isolates rule-based rail closure, and one that isolates
geography-driven road placement.

\subsubsection{Instrument 1: The Larkin Plan Classification}
\label{sssec:iv_larkin}

The first instrument exploits the Larkin Plan (1960--1962), the
government technical study of the rail network described in
Section~\ref{ssec:larkin}. The Plan evaluated segments against an
engineering and freight-traffic standard and designated a subset for
study toward closure; we label these segments \emph{studied} and the
remainder \emph{non-studied}, mapped in
Figure~\ref{fig:larkin_studied}.\footnote{In our cleaned data, studied
segments account for 48.8~percent of network kilometers, compared
with the 39.6~percent the report cites
(Section~\ref{ssec:larkin}). The two figures are consistent with a
definitional difference. The plan's tables assign every studied
segment one of three recommendations
(Section~\ref{ssec:larkin}), and the segments recommended for
\emph{a new study} appear not to be counted as studied in the
report's own accounting. In the raw digitized tables, where the
studied share is 49.0~percent, the new-study segments account for
5{,}197~km; excluding them, the studied share is 38.4~percent of
the approximately 43{,}500-kilometer 1960 network of
Section~\ref{ssec:implementation}. Our instrument retains these
segments: they fell below the freight-density thresholds and were
scrutinized by the plan, which is the classification the instrument
exploits.} The classification predates and is defined
independently of the regional-development outcomes we study. We
construct $\Delta \ln \mathrm{MA}^{\mathrm{LP}}_i$ as the change in log
market access between 1960 and 1986 on a counterfactual network in
which the rail segments removed are exactly those the Larkin Plan
studied, and we use it as an instrument for the realized change
$\Delta \ln \mathrm{MA}^{\mathrm{full}}_i$.

The Larkin classification predicts realized closure: studied segments
were substantially more likely to be closed than non-studied segments
(Section~\ref{ssec:larkin}), so the classification shifts the realized
change $\Delta \ln \mathrm{MA}^{\mathrm{full}}_i$, which is the
relevance condition the instrument must satisfy. The exclusion
restriction requires that the classification be uncorrelated with
district outcome shocks over 1960--1986, conditional on the controls in
equation~(\ref{eq:main}). The assumption is plausible because the Plan
applied a national, traffic-based engineering rule fixed at the start
of the period rather than a forecast of regional growth. It is not
guaranteed: if the traffic rule selected segments in districts already
on a distinct trajectory, studied status would correlate with
pre-existing trends. Sections~\ref{sec:pre_balance}
and~\ref{sec:pre_trends} probe this directly.

\subsubsection{Instrument 2: Hypothetical Road Networks}
\label{sssec:iv_hypo}

The second instrument is a hypothetical road network that connects
Argentina's major cities along the cheapest available routes,
independent of where roads were actually built. The node set is the
cities with at least 15{,}000 inhabitants in 1960 together with all
provincial capitals. We connect these nodes in two ways. Edge cost is
measured either as Euclidean distance or as the least-cost path across
a terrain-based construction-cost surface, on which steeper and more
rugged terrain is more expensive to traverse. From each cost notion we
build a minimum spanning tree --- the lowest-total-cost set of edges
that connects all cities without cycles --- which yields a parsimonious
network that does not depend on which individual city pairs happen to
be linked. This produces four variants: the full bilateral least-cost
network (LCP) and its spanning tree (LCP-MST), and the full Euclidean
network (EUC) and its spanning tree (EUC-MST); Figure~\ref{fig:hypo_networks}
maps the variants. The main specification
uses LCP-MST. We construct $\Delta \ln \mathrm{MA}^{\mathrm{hypo}}_i$
as the change in log market access induced by adding this hypothetical
network.

The hypothetical network depends only on city locations and terrain,
both fixed before the reform and unrelated to district outcome shocks
over 1960--1986. It is relevant to the extent that actual road
construction followed the same low-cost corridors connecting the same
cities; this is an empirical relevance condition, evaluated by the
first-stage fit rather than assumed. It is plausibly excludable because
the least-cost routing reflects engineering geography fixed before the
reform and does not respond to the local economic changes we study. As
with the Larkin instrument, the exclusion restriction is assessed in
the validation package (Sections~\ref{sec:pre_balance}
and~\ref{sec:pre_trends}).

\subsubsection{First Stage}
\label{sssec:first_stage}

The first stage regression is:
\begin{equation}
\Delta \ln \mathrm{MA}^{\mathrm{full}}_i
  = \alpha_1 \cdot \Delta \ln \mathrm{MA}^{\mathrm{LP}}_i
  + \alpha_2 \cdot \Delta \ln \mathrm{MA}^{\mathrm{hypo}}_i
  + X'_i \delta + u_i.
\label{eq:first_stage}
\end{equation}

\subsection{Validation Package}

Both instruments are formula instruments in the sense of
\citet{borusyakhull2023}: each is a market-access index computed from
a counterfactual network, so each combines its source of
identifying variation (the Larkin study designations; the
least-cost-routing geography) with every district's non-random
exposure to that variation through the population weights, district
locations, and the remainder of the network.
\citet{borusyakhull2023} show that instruments of this form can be
correlated with outcome shocks even when the underlying variation is
as good as random, because centrally located districts receive
systematically different instrument values than peripheral ones. The
validation exercises below address this concern directly: the balance
table asks whether the instruments correlate with predetermined
characteristics, the placebo asks whether they predict pre-period
growth, and the main specification conditions on baseline log market
access, baseline log population, and the geographic controls. These
checks partially address, but do not eliminate, the exposure concern.

Sections~\ref{sec:pre_balance} and~\ref{sec:pre_trends} present two
validation exercises that probe the exclusion restriction underlying
the two instruments. Section~\ref{sec:first_stage_tab} reports the
first stage.

A specification with two instruments and one endogenous regressor
imposes one testable restriction: both instruments must identify the
same parameter. \citet{sargan1958} gives the classical test of that
restriction. Standard errors throughout this paper are
heteroskedasticity-robust, so we also compute the
identification-robust counterpart, obtained by minimizing the
\citet{andersonrubin1949} statistic over the parameter and referring
the minimum to a chi-squared distribution with one degree of freedom.
The same statistic delivers Anderson--Rubin confidence sets, which
have correct coverage whatever the strength of the instruments; we
report them where instrument strength is in question.
\citet{andrewsstocksun2019} review both procedures. Rejection of the
restriction does not identify which instrument is at fault. It says
that the two-instrument estimate cannot be read as an estimate of a
common parameter, and that the two instruments should be reported
separately instead. Section~\ref{sec:other_outcomes} is organized
around this test, and Table~\ref{tab:other_outcomes_iv} reports it
for the four outcomes considered there.
% NOTE for Cote/Diego: Table 11 is currently the only table carrying an
% overidentification row. Adding the same row to Tables 9 and 10 is a
% two-line change (sargan_p() from code/analysis/_iv_helpers.R plus an
% add_rows entry). Deliberately not done here: it changes the headline
% tables, so it is a decision for the two of us, not a drafting call.
% Raised as an asymmetry in the PR #141 review.

\subsubsection{Pre-Period Balance}
\label{sec:pre_balance}

Table~\ref{tab:pre_balance} reports regressions of each baseline
district characteristic on the two instruments, partialling out
baseline log market access, baseline log population, and the six
geographic controls of the main specification; a characteristic that
is itself one of these controls is excluded from its own control
set. Each row is a
separate regression of the row-label characteristic; columns report
the instrument coefficients under three specifications: the Larkin
Plan instrument alone (column~1), the hypothetical-road instrument
alone (column~2), and both instruments jointly (column~3). Standard
errors are heteroskedasticity-robust.

The Larkin Plan instrument is uncorrelated with baseline log
population (column~1 coefficient $\balLPPopCoef{}$, $p\balLPPopP{}$).
It shows residual
correlations with baseline urban share ($\balLPUrbshrCoef{}$,
$p\balLPUrbshrP{}$), wheat
suitability ($\balLPWheatCoef{}$, $p\balLPWheatP{}$), pre-1500 caloric potential
($\balLPPreCalCoef{}$, $p\balLPPreCalP{}$), post-1500 caloric potential ($\balLPPostCalCoef{}$,
$p\balLPPostCalP{}$), distance to Buenos Aires ($\balLPDistBACoef{}$, $p\balLPDistBAP{}$), and log
district area ($\balLPAreaCoef{}$, $p\balLPAreaP{}$). These correlations are small in
magnitude but statistically detectable. Wheat suitability, the two
caloric potentials, and distance to Buenos Aires are controls in the
main specification, so their residual
correlations shift from exclusion-restriction threats to a
conditional-independence requirement: the Larkin Plan instrument must
be uncorrelated with the outcome \emph{conditional on} these
geographic controls.

The hypothetical-road instrument shows imbalance at the five-percent
level for \balHypoNImbalanced{} of the nine baseline characteristics
(baseline log population and log district area); the remaining
column~2 coefficients show no correlation at conventional levels.
Baseline log population is a control in the main specification. The
joint specification in
column~3 yields similar results to the single-instrument columns.

\input{../results/tables/table_6_pre_balance.tex}

\subsubsection{Pre-Trends}
\label{sec:pre_trends}

Table~\ref{tab:pre_trends} reports a placebo test. The dependent
variable is the log change in district population between 1947 and
1960 --- a change that predates the 1960--1986 infrastructure
restructuring. The regressor is $\Delta \ln
\mathrm{MA}^{\mathrm{full}}_i$ computed over 1960--1986. If the
infrastructure changes in 1960--1986 are driven by instruments
orthogonal to pre-period district-level dynamics, the coefficient on
post-reform $\Delta \ln \mathrm{MA}^{\mathrm{full}}_i$ in this
placebo regression should be near zero.

The sample is \placeboNIVBoth{} districts --- the subset for which the 1947 census
provides comparable population counts under the time-invariant
district boundaries, as discussed in Section~\ref{sec:data}. Columns~(1) to~(4)
match the structure of Tables~\ref{tab:population_iv}
and~\ref{tab:sectoral_iv}: OLS, IV-LP, IV-Hypo, IV-Both.

The controls in this table differ from those in
Tables~\ref{tab:pre_balance}, \ref{tab:first_stage}
and~\ref{tab:population_iv} in one term: the population baseline is
the 1947 level, not the 1960 level. The reason is that the outcome
here is $\ln \mathrm{pop}_{1960} - \ln \mathrm{pop}_{1947}$, so the
1960 level is the terminal value of the window under test, and
conditioning on it would mean conditioning on a component of the
outcome. The 1947 level is the initial value, a convergence control.
The baseline market-access control is retained.
Table~\ref{tab:placebo_ladder} reports the estimate under four
baseline control sets, including the set that drops the
market-access control; the estimate is sensitive to that choice, and
Section~\ref{sec:limitations} returns to it.

The OLS coefficient is \placeboOLSCoef{} (\placeboOLSSE{}), with
$p=\placeboOLSP{}$. The IV-LP estimate is \placeboIVLPCoef{} (\placeboIVLPSE{}), the IV-Hypo estimate is
\placeboIVHCoef{} (\placeboIVHSE{}), and the combined IV estimate is \placeboCoefIVBoth{} (\placeboSEIVBoth{}),
with $p=\placeboIVBP{}$. The first-stage $F$ for the
IV-Hypo column on this sample of \placeboNIVBoth{} districts is \placeboFHypo{}, below
conventional weak-instrument thresholds; the IV-Hypo column is
therefore not interpretable in this specification. The first-stage
$F$ for IV-LP is \placeboFLP{} and for IV-Both is \placeboFBoth{}.

The combined-IV point estimate for the placebo outcome is positive
and rejects a zero coefficient at the ten-percent level, though not
at five percent; the OLS estimate has the same sign and does not
reject. This is not a clean null. Two readings are
consistent with the evidence. First, there may be a genuine
pre-existing trend correlated with the eventual infrastructure
changes of 1960--1986: districts that were growing fastest in
1947--1960 may have been more likely to gain (or lose) market access
in the later period. Second, the placebo outcome is defined on a
subsample ($N=\placeboNIVBoth{}$) that is non-random with respect to the full
estimation sample ($N=\mainPopNIVBoth{}$), and the correlations documented in
Section~\ref{sec:pre_balance} for the Larkin Plan instrument may
interact with that subsample selection.

Section~\ref{sec:results_robustness} reports the main-spec population
regression refit on the same \placeboNIVBoth{}-district subsample. On that
subsample, the IV-Both coefficient on $\Delta \ln
\mathrm{MA}^{\mathrm{full}}_i$ is \subIVBCoef{} (\subIVBSE{}), with $p=\subIVBP{}$. This
coefficient is similar in magnitude to the placebo coefficient of
\placeboCoefIVBoth{} reported above. The two readings offered in the preceding
paragraph are therefore not distinguishable by this test: the pre-
reform pattern on the \placeboNIVBoth{} subset is of the same order as the
post-reform main-spec coefficient on the same subset. We retain the
main specification on the full sample as the headline estimate and
report the placebo-subsample result in the robustness panel.

\input{../results/tables/table_7_pre_trends.tex}

\subsubsection{First-Stage Results}
\label{sec:first_stage_tab}

Table~\ref{tab:first_stage} reports the first-stage regressions.
Results are described in Section~\ref{sec:first_stage}.

\input{../results/tables/table_8_first_stage.tex}

\subsection{Additional Robustness Approaches}

We assess robustness by varying the distance decay parameter
($\theta$), the set of controls, the sample (excluding major cities,
border regions), and the time period (1960--1970, 1970--1991).
Results are reported in Table~\ref{tab:robustness}.
